\section{Cost model and metrics}
\label{sec:cost}

We score a request by what it costs to get the pod to finish, counting failed attempts, retries and lateness,
all simulated from the run's measurements. Each axis fails the way it fails on Kubernetes. If memory is too
low, the container is OOM-killed when the run reaches its peak (at the measured fraction of its runtime) and is
retried with double the memory, for up to eight attempts. If CPU is too low, the run is throttled: each second
that wanted $u$ cores takes $\max(1, u/\text{limit})$ seconds, replayed against the measured per-second series.
If the GPU is too small, the run hits a CUDA out-of-memory error and is retried on the next machine up, for up
to four attempts. A GPU task given no GPU fails at start and is retried on a T4. A CPU task given a GPU runs
normally and pays the GPU node's price.

A CPU pod is priced by its dominant share of an \code{m6i.xlarge} (4 vCPU, 16\,GiB, \$0.192 per hour), that is
$p(m,c) = 0.192\cdot\max(c/4,\ \text{GiB}(m)/16)$, and a GPU pod pays the on-demand price of its machine. Every
attempt pays a start-up overhead: 0.02 hours on a CPU node and 0.05 hours on a GPU node, for image pull, device
initialization and reading inputs. We normalize every run to one nominal hour of work, so that short measured
runs stand in for longer production tasks. A run that exhausts its attempts pays for the attempts it made; this
happens for fewer than 0.4\% of evaluated pods for any approach.

Two penalties price time. Each retry adds the cost of running the label request once, which stands for the
delay of running again, and each hour of lateness from throttling adds the same amount per hour. The label
request's cost is a fixed quantity for each run, so the penalties do not depend on the request being scored.
These defaults match an orchestrator that retries automatically. A raw Job with \code{backoffLimit: 0} fails and
waits for someone to resubmit it, so \cref{sec:res-sens} also scores every approach with retry penalties of 0,
5 and 20. Because costs include these penalties, the ``dollars'' in this paper are priced dollars, not a cloud
bill.

\paragraph{The cheapest request.}
For each pod, the benchmark searches the grid for the request with the lowest cost under this model: every
memory and CPU cell with no GPU for a CPU task (272 requests), and every GPU cell for a GPU task. This is the
denominator of regret, and no request can beat it. It is often cheaper than the label, the smallest cell that
fits: asking for one CPU cell below the busiest 15 seconds only throttles that window, and the cores saved over
the rest of the run can cost less than the throttling. Against the label, our first release reported negative
regret for some approaches on RightSizeCPU; against the cheapest request, regret is never negative.

\paragraph{Metrics.}
The headline metric is \emph{regret}, the mean over pods of an approach's cost divided by the cheapest
request's cost, minus one. Every pod counts equally however much it costs, and this is the quantity every
approach tunes on validation. \emph{Dollar regret} is total cost over total cheapest cost, minus one, so large
jobs count for more. We also report the share of runs whose first attempt succeeds. For pairs of pods that come
from the same run but differ in one detail of the pod spec, we check whether an approach's memory request moves
in the same direction as the label (\cref{sec:res-manifest}), separately for pairs whose label moved and pairs
whose label did not.

\paragraph{Uncertainty.}
Runs of the same task share code, and a run's Job and Flyte pods share a workload, so pods are not independent.
Intervals are 95\% percentile intervals from a clustered Poisson bootstrap~\citep{efron1993bootstrap,field2007clustered}
with 2,000 replicates. By default each task, with all of its runs and both pods of each run, gets one Poisson(1)
weight; this measures uncertainty about new tasks from the same families. A claim about new kinds of code has
only seven held-out families behind it, so for the unseen split we also cluster by family, which gives much
wider intervals. Differences between approaches and between pod types are bootstrapped as paired differences on
the same pods. Variation between trainings is reported separately, across three seeds (\cref{sec:res-seeds}).
